\documentclass[12pt]{article}
\usepackage[a4paper,margin=1in]{geometry}
\usepackage{amsmath,amssymb,amsthm}

\title{A Simple Derivation from a Block Intuition to a Periodic Congruence Formula}
\author{Matteo Calvigioni}
\date{}

\begin{document}

\maketitle

\noindent
\textbf{Premise.}
The starting point was not a theorem, but only a simple \emph{block intuition}. 
One may look at an object of the form
\[
L\,n,
\]
where \(L\) is a letter and \(n\) is a decimal block of digits, and associate to the letter a positive integer \(p\).
From this informal point of view, one is led to consider a decomposition into three parts:
a left block \(p\), a middle block \(n\), and a right block \(n \bmod p\).
This suggests the identity
\[
2S(n)=p+(n \bmod p),
\]
where \(S(n)\) denotes the sum of the decimal digits of \(n\).

\medskip

\noindent
From here, one may forget the original block picture and study the relation on its own.

\section*{Step 1: From the identity to the congruence}

Assume
\[
2S(n)=p+(n \bmod p).
\]
Reducing both sides modulo \(p\), we get
\[
2S(n)\equiv p+(n \bmod p)\pmod p.
\]
Since
\[
p\equiv 0 \pmod p
\]
and, by definition,
\[
n \bmod p \equiv n \pmod p,
\]
it follows that
\[
2S(n)\equiv n \pmod p.
\]

Thus the original identity implies the congruence
\[
2S(n)\equiv n \pmod p.
\]

\section*{Step 2: Explicit form of the congruence}

By the definition of congruence modulo \(p\),
\[
2S(n)\equiv n \pmod p
\]
is equivalent to the existence of some integer \(k\in\mathbb{Z}\) such that
\[
2S(n)-n=pk.
\]
This is only the explicit form of the same congruence.

\section*{Step 3: Decimal expansion}

From now on, we assume that we are working in base \(10\).

Let
\[
n=d_0+d_1 10+d_2 10^2+\cdots+d_m 10^m,
\]
where each digit satisfies
\[
d_i\in\{0,1,\dots,9\}.
\]
Then the sum of the digits is
\[
S(n)=d_0+d_1+\cdots+d_m.
\]

Substituting this into the congruence
\[
2S(n)\equiv n \pmod p,
\]
we obtain
\[
2(d_0+d_1+\cdots+d_m)\equiv d_0+d_1 10+d_2 10^2+\cdots+d_m 10^m \pmod p.
\]
Bringing everything to the left-hand side gives
\[
\sum_{i=0}^m d_i(2-10^i)\equiv 0 \pmod p.
\]
Equivalently,
\[
2S(n)-n=\sum_{i=0}^m d_i(2-10^i).
\]

So the original congruence may be rewritten as
\[
\sum_{i=0}^m d_i(2-10^i)\equiv 0 \pmod p.
\]

\section*{Step 4: Periodicity of \(10^i \pmod p\)}

Now assume
\[
\gcd(10,p)=1.
\]
Then \(10\) is invertible modulo \(p\). Hence the powers
\[
1,10,10^2,10^3,\dots
\]
belong to a finite multiplicative system modulo \(p\), so two of them must eventually coincide.
That is, there exist integers \(r>s\geq 0\) such that
\[
10^r\equiv 10^s \pmod p.
\]
Since \(10^s\) is invertible modulo \(p\), we may cancel it and obtain
\[
10^{r-s}\equiv 1 \pmod p.
\]
Therefore there exists a least positive integer \(t\) such that
\[
10^t\equiv 1 \pmod p.
\]
This integer \(t\) is the period of the sequence
\[
10^i \pmod p,
\]
in the sense that
\[
10^{i+t}\equiv 10^i \pmod p
\qquad\text{for all } i\geq 0.
\]

\section*{Step 5: Periodicity of \(2-10^i \pmod p\)}

Since
\[
10^{i+t}\equiv 10^i \pmod p,
\]
we immediately get
\[
2-10^{i+t}\equiv 2-10^i \pmod p
\qquad\text{for all } i\geq 0.
\]
So the sequence
\[
(2-10^i)\pmod p
\]
is also periodic (with period dividing \(t\), and certainly with period \(t\)).

\section*{Step 6: Grouping the digits by residue classes modulo \(t\)}

Because the coefficients \(2-10^i\) repeat modulo \(t\), one may group the indices according to their residue classes modulo \(t\).

For each \(r\in\{0,1,\dots,t-1\}\), define
\[
C_r(n):=\sum_{\substack{0\leq i\leq m\\ i\equiv r \;(\mathrm{mod}\; t)}} d_i.
\]
This is the sum of the digits of \(n\) whose positions are congruent to \(r\) modulo \(t\).

Since
\[
10^i\equiv 10^r \pmod p \qquad \text{whenever } i\equiv r \pmod t,
\]
we also have
\[
2-10^i\equiv 2-10^r \pmod p.
\]
Hence
\[
\sum_{i=0}^m d_i(2-10^i)
\equiv
\sum_{r=0}^{t-1}(2-10^r)\!\!\sum_{\substack{0\leq i\leq m\\ i\equiv r\;(\mathrm{mod}\; t)}} d_i
\pmod p.
\]
By the definition of \(C_r(n)\), this becomes
\[
\sum_{i=0}^m d_i(2-10^i)
\equiv
\sum_{r=0}^{t-1}(2-10^r)C_r(n)
\pmod p.
\]

Since the original congruence is equivalent to
\[
\sum_{i=0}^m d_i(2-10^i)\equiv 0 \pmod p,
\]
we finally obtain the formula
\[
\boxed{
\sum_{r=0}^{t-1}(2-10^r)C_r(n)\equiv 0 \pmod p
}
\]
where \(t\) is the period of \(10^i \pmod p\), under the assumption \(\gcd(10,p)=1\).

\section*{Conclusion}

Starting from an informal block intuition, one is naturally led to the identity
\[
2S(n)=p+(n \bmod p),
\]
then to the congruence
\[
2S(n)\equiv n \pmod p,
\]
then to its explicit form
\[
2S(n)-n=pk,
\]
and finally, after writing \(n\) in decimal notation, to
\[
\sum_{i=0}^m d_i(2-10^i)\equiv 0 \pmod p.
\]
Under the additional assumption \(\gcd(10,p)=1\), the periodicity of \(10^i \pmod p\) allows us to group the digits by residue classes modulo the period \(t\), obtaining the final formula
\[
\sum_{r=0}^{t-1}(2-10^r)C_r(n)\equiv 0 \pmod p.
\]

\medskip

\noindent
This argument is elementary, but it already shows that the original relation is controlled by a periodic structure coming from the powers of \(10\) modulo \(p\).
\section*{Remarks}

The congruence
\[
2S(n)\equiv n \pmod{p}
\]
reveals two distinct structural regimes, depending on the arithmetic properties of $p$.

\medskip

\textbf{(1) The coprime case: $\gcd(10,p)=1$.}

In this case, the sequence $(10^i \bmod p)$ is periodic. Let $t$ be the multiplicative order of $10$ modulo $p$, so that
\[
10^t \equiv 1 \pmod{p}.
\]
Then the coefficients $2 - 10^i$ are also periodic with period $t$, and the congruence
\[
\sum_{i=0}^m d_i(2-10^i)\equiv 0 \pmod{p}
\]
can be grouped according to residue classes modulo $t$.

This leads naturally to the quantities
\[
C_r(n) := \sum_{\substack{i \equiv r \ (\mathrm{mod}\ t)}} d_i,
\]
and the condition becomes a linear relation among the $C_r(n)$.

\medskip

\textbf{(2) The non-coprime case: $\gcd(10,p)\neq 1$.}

If $p$ shares a factor with $10$, the behavior changes drastically. In particular, powers of $10$ may eventually vanish modulo $p$.

For example, if $p=25$, then
\[
10^2 \equiv 0 \pmod{25},
\]
and therefore
\[
10^i \equiv 0 \pmod{25} \quad \text{for all } i \ge 2.
\]
In this situation, the congruence simplifies to
\[
d_0 + 17d_1 + 2\sum_{i=2}^m d_i \equiv 0 \pmod{25},
\]
so the condition depends only on the last two digits and on the sum of the remaining digits.

\medskip

\textbf{(3) A stronger identity in special cases.}

In some instances, one observes the stronger relation
\[
2S(n) = p + (n \bmod p),
\]
which implies the congruence. This does not hold in general, but it provides a useful heuristic starting point.

\medskip

The example $n=57342$ for $p=25$ illustrates this behavior, showing how a simple digit-based condition can encode a modular relation.


\end{document}
